\textbf{2.\ \textsc{Gas} (6 points)} --- \textit{Jaan Kalda.}

A box of volume $V$ contains $\nu$ moles of a monoatomic gas of molar mass $\mu$
at a negligibly low temperature. The box stops instantaneously after having
moved with a constant speed $v$ (much greater than the thermal speed).

\textbf{i)} \textit{(2 points)} Find the temperature inside the box after
thermalization.

\textbf{ii)} \textit{(2 points)} Find pressure to the front wall (the wall that
was at the front while the box was moving) immediately after stopping.

\textbf{iii)} \textit{(2 points)} Now, a spherical ball of gas (helium,
$\mu = 4\,\mathrm{g/mol}$ of radius $r = 1\,\mathrm{cm}$, at temperature
$T = 300\,\mathrm{K}$ is surrounded by vacuum. The mean free path of the
molecules is much larger than $r$. At a certain moment, the walls of the ball
break and after $\tau = 5\,\mathrm{ms}$, a part of the gas is trapped by erecting
instantaneously walls forming a cubical vessel of volume $V = 1\,\mathrm{m^3}$.
Find the temperature of the gas inside the cube $T$' after thermalization;
neglect the thermal capacitance of the walls of the cubical vessel.
$R = 8.31\,\mathrm{J/kg\,K}$.
